%%%%%%%%%%%%%%%%%%%%%%%%%%%%%%%%%%%%%%%%%%%%%%%%%%%%%%%%%%%%%%%%%%%%%%%%%%%%
%%%%%%%%%%%%%%%%%%%%%%%%%%%%%%%%%%%%%%%%%%%%%%%%%%%%%%%%%%%%%%%%%%%%%%%%%%%%
%%%%%%%%%%%%%%%%%%%%%%%%%%%%%%%%%%%%%%%%%%%%%%%%%%%%%%%%%%%%%%%%%%%%%%%%%%%%

% TEXT FORMATTING

\newcommand{\ie}{i.\,e.\ }
\newcommand{\Ie}{I.\,e.\ }
\newcommand{\eg}{e.\,g.\ }
\newcommand{\Eg}{E.\,g.\ }
\newcommand{\confer}{cf.\ }
\renewcommand{\quote}[1]{``#1''}
\newcommand{\emphquote}[1]{\emph{\quote{#1}}}

\newcommand{\assumptionPreWord}{\textcolor{red}{assumption(!)}\,}

%%%%%%%%%%%%%%%%%%%%%%%%%%%%%%%%%%%%%%%%%%%%%%%%%%%%%%%%%%%%%%%%%%%%%%%%%%%%

% GENERAL MATH MODE FORUMALS
\newcommand{\ldef}{\coloneqq}
\newcommand{\rdef}{\eqqcolon}
\newcommand{\setseparator}{\,\vert\,}
\newcommand{\powerset}[1]{{\mathcal P(#1)}}

\newcommand{\N}{{\mathbb{N}}}
\newcommand{\A}{{\mathbb{A}}}
%\renewcommand{\P}{\mathbb{P}}
%\newcommand{\F}{\mathbb{F}}
\newcommand{\Z}{{\mathbb{Z}}}
\newcommand{\Q}{{\mathbb{Q}}}
\newcommand{\R}{{\mathbb{R}}}
\newcommand{\C}{{\mathbb{C}}}
\newcommand{\adicsQ}[1]{{\Q_{#1}}}

%%%%%%%%%%%%%%%%%%%%%%%%%%%%%%%%%%%%%%%%%%%%%%%%%%%%%%%%%%%%%%%%%%%%%%%%%%%%

% UNSORTED COMMANDS
\newcommand{\yoneda}{{\mathcal Y}}
\newcommand{\iso}{\cong}
\renewcommand{\equiv}{\simeq}
\newcommand{\grothendieckaxiom}[1]{(AB#1)}

% CATEGORIES
%\newcommand{\Ob}{\operatorname{Ob}}
%\newcommand{\Mor}{\operatorname{Mor}}
%\newcommand{\End}{\operatorname{End}}
%\newcommand{\Aut}{\operatorname{Aut}}


% general commands regarding category theory
\DeclareMathAlphabet\mathbfcal{OMS}{cmsy}{b}{n}

\newcommand{\category}[1]{{\mathbfcal{#1}}}
%\newcommand{\categoryname}[1]{{\normalfont{\textbf{#1}}}}
\newcommand{\categoryname}[1]{{\mathbf{#1}}}

\newcommand{\Hom}{{\operatorname{Hom}}}
\newcommand{\eHom}{{\underline{\operatorname{Hom}}}}
\newcommand{\intHom}{{\operatorname{\mathscr{H}\!om}}}
\newcommand{\structHom}[1]{\Hom_{\struct{#1}}}
\newcommand{\eStructHom}[1]{\eHom_{\struct{#1}}}
\newcommand{\eRHom}{\RightD\eHom}
\newcommand{\intRHom}{\RightD\intHom}
\newcommand{\eStructRHom}[1]{\eRHom_{\struct{#1}}}

\newcommand{\id}{\mathrm{id}}
\newcommand{\shortisorightarrow}{\stackrel{\sim\;\;}{\smash{\to}\rule{0pt}{0.2ex}}}
\newcommand{\isorightarrow}{\stackrel{\sim\;\;\;\;\;}{\smash{\longrightarrow}\rule{0pt}{0.2ex}}}
\newcommand{\shortisoleftarrow}{\stackrel{\;\;\sim\,}{\smash{\leftarrow}\rule{0pt}{0.2ex}}}
\newcommand{\isoleftarrow}{\stackrel{\;\;\;\;\,\sim\,}{\smash{\longleftarrow}\rule{0pt}{0.2ex}}}


% CATEGORY SCHEMAS
\newcommand{\op}{{\mathrm{op}}}
\newcommand{\Overcategory}[2]{{#1\categoryname{/}#2}}
\newcommand{\functorCategory}[2]{{\categoryname{Func}\mleft(#1, #2\mright)}}

\newcommand{\preSheavesSymbol}{{\categoryname{PSh}}}
\newcommand{\sheavesSymbol}{{\categoryname{Sh}}}
\newcommand{\Presheaves}[2]{{\preSheavesSymbol{\mleft(#1, #2\mright)}}}
\newcommand{\Sheaves}[2]{{\sheavesSymbol{\mleft(#1, #2\mright)}}}

\newcommand{\smallPresheaves}[2]{{\categoryname{P\underline{s}h}{\mleft(#1, #2\mright)}}}
\newcommand{\smallSheaves}[2]{{\categoryname{\underline{s}h}{\mleft(#1, #2\mright)}}}

\newcommand{\ProCategory}[1]{{\categoryname{Pro}{\mleft(#1\mright)}}}

\newcommand{\ladj}{\dashv}
\newcommand{\unit}{\eta}
\newcommand{\coUnit}{\epsilon}

\newcommand{\longxrightarrow}[1]{\xrightarrow{\;#1\;}}
\newcommand{\xleftrightarrows}[2]{\xleftrightharpoons[\;#2\;]{\;#1\;}}

% SPECIFIC CATEGORIES
\newcommand{\Set}{{\categoryname{\underline{S}et}}}
\newcommand{\boundedSet}[1]{{\Set_{#1}}}
\newcommand{\set}{{\categoryname{\underline{s}et}}}
\newcommand{\Top}{{\categoryname{Top}}}
\newcommand{\Schemes}{{\categoryname{Schemes}}}
\newcommand{\Ab}{{\categoryname{Ab}}}
\newcommand{\TopAb}{{\categoryname{TopAb}}}
\newcommand{\dmod}[1]{{\prescript{\delta}{#1}{\categoryname{Mod}}}}
\newcommand{\commutativeRings}{\categoryname{CR1ng}}

% LIMITS & COLIMITS
\newcommand{\ConstantFunctor}[1]{{{\mathbf 1}_{#1}}}
\newcommand{\limit}{\varprojlim}
\newcommand{\colimit}{\varinjlim}
\newcommand{\homotopyLimit}{\mathrm{h}\!\limit\nolimits}
\newcommand{\homotopyColimit}{\mathrm{h}\!\colimit\nolimits}

\DeclareMathOperator{\equalizer}{Eq}
%\DeclareMathOperator{\CoEq}{CoEq}
\DeclareMathOperator{\kernel}{Ker}
\DeclareMathOperator{\coKernel}{CoKer}
\DeclareMathOperator{\image}{Img}
%\DeclareMathOperator{\CoImage}{CoImg}

% SITES AND SITE SPECIFIC COMMANDS
\newcommand{\site}[1]{{\categoryname{#1}}}

\newcommand{\covers}[1]{{\operatorname{Cov}(#1)}}

\newcommand{\open}[1]{{\categoryname{Ouv}(#1)}}
\newcommand{\etalesite}[1]{{{#1}_{\text{ét}}}}
\newcommand{\proetalesite}[1]{{{#1}_{\text{pro-ét}}}}

\newcommand{\CHaus}{{\categoryname{CHaus}}}
\newcommand{\boundedCHaus}[1]{{\CHaus_{#1}}}

\newcommand{\proFiniteSets}{\categoryname{ProFin}}
\newcommand{\boundedProFiniteSets}[1]{\proFiniteSets_{#1}}

\newcommand{\extremallyDisconnectedSets}{\categoryname{ExDisc}}
\newcommand{\boundedExtremallyDisconnectedSets}[1]{\extremallyDisconnectedSets_{#1}}

% CONDENSED
\newcommand{\condensedSymbol}{{\categoryname{Cond}}}
\newcommand{\boundedCondensed}[2]{{\condensedSymbol_{#1}(#2)}}
\newcommand{\boundedCondensedSets}[1]{{\condensedSymbol\Set_{#1}}}
\newcommand{\associated}[1]{\underline{#1}}
\newcommand{\condensed}[1]{{\condensedSymbol(#1)}}
\newcommand{\condensedSets}{{\condensedSymbol\Set}}
\newcommand{\condensedAb}{{\condensedSymbol\Ab}}
\newcommand{\topologizedUnderlying}[1]{{#1(\ast)_\mathrm{top}}}
\newcommand{\resTo}[1]{{\vert_{#1}}}
\newcommand{\enlargementsymbol}[2]{{#1\rightsquigarrow #2}}
\newcommand{\CGWH}{C.G.W.H. }

% SHEAF SPECIFIC COMMANDS
\newcommand{\sheaf}[1]{{\mathcal{#1}}}
\newcommand{\sheafificationsymbol}{\sharp}
\newcommand{\sheafification}[1]{{\mleft(#1\mright)^\sheafificationsymbol}}
\renewcommand{\restriction}{\vert}

% CARDINALS & ORDINALS
\newcommand{\Ordinals}{{\categoryname{Ord}}}
\newcommand{\Cardinals}{{\categoryname{Card}}}
\newcommand{\cofinality}[1]{{\operatorname{cof}(#1)}}

% GENERAL TOPOLOGY
\newcommand{\topology}[1]{{\mathscr{#1}}}
\newcommand{\boundedkSpaces}[1]{{\mathrm{k}\Top^{#1}}}
\newcommand{\kSpaces}{{\mathrm{k}\Top}}
\newcommand{\kificationsymbol}{{\mathrm{cg}}}
\newcommand{\boundedkificationsymbol}[1]{{#1\mathrm{-cg}}}
\newcommand{\stonecechsymbol}{{\beta}}
\newcommand{\stonecech}[1]{{\stonecechsymbol #1}}
\newcommand{\filterkernel}{{\operatorname{ker}}}
\newcommand{\tOne}{{T_1}}
\newcommand{\tOneSpaces}{{\Top_1}}

% RINGS AND MODULES
\newcommand{\abelianObject}[1]{\Ab\mleft(#1\mright)}
\newcommand{\structsymbol}{{\mathcal{O}}}
\newcommand{\struct}[1]{{\structsymbol_{#1}}}
\newcommand{\stalkStruct}[2]{{\structsymbol_{#1, #2}}}

\newcommand{\structTensor}[1]{\otimes_{\struct{#1}}}
\newcommand{\associatedSheaf}[1]{\widetilde{#1}}

\newcommand{\constantSheaf}[1]{{\underline{#1}}}
\newcommand{\free}[2]{{#1[#2]}}
\newcommand{\freeAbelian}[1]{{\free {\associated{\Z}} {#1}}}

\newcommand{\scalarRestriction}[1]{{#1^\ast}}
\newcommand{\scalarExtension}[1]{{#1_{(!)}}}

\newcommand{\ringedSpaces}{{\categoryname{RingSp}}}
\newcommand{\locallyRingedSpaces}{{\categoryname{LocRingSp}}}
\newcommand{\ringedSites}{{\categoryname{RingSit}}}
\newcommand{\pMod}[1]{{\categoryname{pMod}(#1)}}
\renewcommand{\mod}[1]{{\categoryname{Mod}(#1)}}
\newcommand{\maybePMod}[1]{{\categoryname{(p)Mod}(#1)}}
\newcommand{\structPMod}[1]{{\pMod{\struct{#1}}}}
\newcommand{\structMod}[1]{{\mod{\struct{#1}}}}
\newcommand{\qCoh}[1]{{\categoryname{qCoh}(\struct{#1})}}
\newcommand{\Coh}[1]{{\categoryname{Coh}(\struct{#1})}}

% SCHEMES AND ADIC SPACES
\newcommand{\ideal}[1]{{\langle #1\rangle}}
\newcommand{\Spec}{\operatorname{Spec}}
\newcommand{\Spa}{\operatorname{Spa}}
\newcommand{\ad}{{\mathrm{ad}}}
\newcommand{\quot}[1]{{\operatorname{Quot}\mleft(#1\mright)}}
%\newcommand{\AffSchemesOver}[1]{{\catname{AffSchemes}_{#1}}}


% ANALYTIC RINGS
\newcommand{\preAnaRing}[1]{{\mathfrak{#1}}}
\newcommand{\underlying}[1]{{\underline{\preAnaRing{#1}}}}
\newcommand{\ufree}[2]{{\free {\underlying{#1}} {#2}}}
\newcommand{\pfree}[2]{{\free {\preAnaRing{#1}} {#2}}}

\newcommand{\preAnalyticRings}{\categoryname{PAnaRing}}
\newcommand{\analyticRings}{\categoryname{AnaRing}}
\newcommand{\presentableMod}[1]{{\categoryname{Pres}(#1)}}
\newcommand{\pPresentableMod}[1]{{\categoryname{Pres}(\preAnaRing{#1})}}
\newcommand{\completeMod}[1]{{\categoryname{Comp}(#1)}}
\newcommand{\pCompleteMod}[1]{{\categoryname{Comp}(\preAnaRing{#1})}}

\newcommand{\completionSymbol}[1]{#1}
\newcommand{\derivedCompletionSymbol}[1]{{\LeftD#1}}
\newcommand{\completion}[2]{\mleft(#1\mright)^{\completionSymbol{#2}}}
\newcommand{\completionAsTensor}[3]{#1 \otimes_{#2} #3}
\newcommand{\derivedCompletion}[2]{\mleft(#1\mright)^{\derivedCompletionSymbol{#2}}}
\newcommand{\derivedCompletionAsTensor}[3]{#1 \Dotimes_{#2} #3}
\newcommand{\completedScalarExtension}[2]{{#1_{(!)}^{#2}}}
\newcommand{\derivedCompletedScalarExtension}[2]{{\LeftD#1_{(!)}^{#2}}}

\newcommand{\continuous}[2]{{C(#1, #2)}}
\newcommand{\continuousZero}[2]{{C_0(#1, #2)}}
\newcommand{\measures}[2]{{\mathcal{M}(#1, #2)}}
\renewcommand{\d}{\,\mathrm{d}}
\newcommand{\laurent}[2]{{#1(\!(#2)\!)}}
\newcommand{\series}[2]{{#1[\![#2]\!]}}


% DERIVED ALGEBRA
\newcommand{\derived}[1]{{\categoryname{D}(#1)}}
\newcommand{\derivedPlus}[1]{{\categoryname{D}^+(#1)}}
\newcommand{\derivedMinus}[1]{{\categoryname{D}^-(#1)}}
\newcommand{\derivedFlat}[1]{{\categoryname{D}^\flat(#1)}}
%\newcommand{\Injectives}{\mathcal{I}}
%\newcommand{\Projectives}{\mathcal{P}}
\newcommand{\RightD}{\mathcal{R}}
\newcommand{\LeftD}{\mathscr{L}}
\newcommand{\leftKanExtension}{\operatorname{Lan}}
%\newcommand{\RAN}{\operatorname{Ran}}

\newcommand{\Ext}{{\operatorname{Ext}}}
\newcommand{\eExt}{{\operatorname{\underline{Ext}}}}
\newcommand{\intExt}{{\operatorname{\mathscr{E}xt}}}
\newcommand{\Dotimes}{\otimes^\LeftD }


\newcommand{\simplicial}[1]{{\categoryname{Simp}(#1)}}
\newcommand{\simplicialSets}{{\categoryname{s}\Set}}
\newcommand{\largeSimplicialSets}{{\categoryname{s\underline{\mathfrak{S}}et}}}
\newcommand{\horn}[2]{{\Lambda^{#1}_{#2}}}

\DeclareMathOperator{\totalComplex}{Tot}

\newcommand{\nerveSymbol}{{\mathcal N}}
\newcommand{\nerve}[1]{{\nerveSymbol(#1)}}
\newcommand{\homotopyCategory}[1]{{h(#1)}}
\newcommand{\categoryOfSmallCategories}{{\categoryname{Cat}}}
\newcommand{\categoryOfCategories}{{\categoryname{CAT}}}
\newcommand{\categoryOfInftyCategories}{{\categoryname{CAT}_\infty}}
\newcommand{\constantFunctor}[1]{{\mathbf{1}_{#1}}}
\newcommand{\mappingSpace}{{\operatorname{Map}}}
\newcommand{\complexes}[1]{{\categoryname{Ch}(#1)}}
\newcommand{\connectiveComplexes}[1]{{\categoryname{Ch}_{\ge 0}(#1)}}

\newcommand{\blownUpSimplex}{{\mathfrak{C}}}
\newcommand{\simpliciallyEnrichedSmallCategories}{\categoryOfSmallCategories_\simplicialSets}
\newcommand{\simpliciallyEnrichedCategories}{{\categoryOfCategories_\simplicialSets}}
\newcommand{\homotopyCoherentNerve}[1]{{\nerveSymbol_\mathrm{hcoh}(#1)}}

\newcommand{\sTruncation}[1]{{\sigma_{#1}}}
\newcommand{\cTruncation}[1]{{\tau_{#1}}}
\newcommand{\homotopyCategoryOfComplexes}[1]{{K(#1)}}

\newcommand{\inftyD}{{\mathcal D}}

% DUALITY
\newcommand{\dual}{\vee}
\newcommand{\exterior}{\bigwedge\nolimits}
\newcommand{\cotangent}{\Omega}
\newcommand{\canonical}{\omega}
\newcommand{\Trace}{\int }
%\newcommand{\codim}{\operatorname{codim}}
%\newcommand{\projdim}{\operatorname{pd}}
%\newcommand{\depth}{\operatorname{depth}}

\newcommand{\support}[1]{{\operatorname{supp}(#1)}}
\newcommand{\compactlySupportedGlobalSections}{\Gamma_{\mathrm c}}
\newcommand{\compactlySupportedCohomology}{H_{\mathrm c}}

%%%%%%%%%%%%%%%%%%%%%%%%%%%%%%%%%%%%%%%%%%%%%%%%%%%%%%%%%%%%%%%%%%%%%%%%%%%%
%%%%%%%%%%%%%%%%%%%%%%%%%%%%%%%%%%%%%%%%%%%%%%%%%%%%%%%%%%%%%%%%%%%%%%%%%%%%
%%%%%%%%%%%%%%%%%%%%%%%%%%%%%%%%%%%%%%%%%%%%%%%%%%%%%%%%%%%%%%%%%%%%%%%%%%%%